% 原文 pp.237--241
% 原文 p.237

\subsectionlink{まとめ}
デコヒーレンス模型についての本章は、やや長く、技術的になった。そこで主要な結果を、いくつかまとめよう。

% 原文 p.238
\hypertarget{srcpage:238}{}\label{srcpage:238}
\subsubsectionlink*{標準模型への写像}
関心のある物理系の広いクラスは、スピン$1/2$粒子（状態空間が実効的に異なる二つの基底状態に縮約される系の場合）、または、位置と運動量という連続的な正準座標で記述され、あるポテンシャルの中を運動する粒子の、いずれかで表せる。

環境は、非局在化したボソン場モードの準連続体に対応する調和振動子の集まりとして、あるいはスピンで表す局在モードの集まりとして、模型化できる。スピン環境は、中心系に強く結合する低温環境の、適切な模型となる場合が多い。

\subsubsectionlink*{量子Brown運動}
これは、中心粒子が実空間のあるポテンシャルの中を運動し、互いに相互作用しない調和振動子の環境と双線形に相互作用する模型である。中心系を調和振動子で表す場合に注目した。

この模型のBorn--Markovマスター方程式を導いた［(5.25)参照］。
\chfiveeq{202}
右辺の第1項は、$\widetilde\Omega^2$で定量化する振動数シフトを伴うユニタリーな運動を記述する。第2項は、$\gamma$に比例する率での運動量減衰に対応する。第3項は、距離$\Delta X$にわたる空間的コヒーレンスの、率$D(\Delta X)^2$でのデコヒーレンスを記述する。第4項もデコヒーレンスを表すが、特に高温では、しばしば無視できる。マスター方程式(5.202)に現れるすべての係数は、環境の物理的特徴をまとめた、環境スペクトル密度に強く依存する。

高温極限では、多くの場合、中心系の縮約された運動は、Caldeira--Leggettマスター方程式として知られる、(5.202)を簡単化した形でよく近似される。その形は［(5.57)参照］
\chfiveeq{203}
である。ここでも右辺の第1項は、環境によって繰り込まれた自己Hamiltonian $\hat H'_{\mathrm S}$が支配するユニタリーな発展を表し、第2項は散逸を記述し、第3項は、浴の温度に線形に増加する率での、空間的重ね合わせのデコヒーレンスに対応する。

% 原文 p.239
\hypertarget{srcpage:239}{}\label{srcpage:239}
Caldeira--Leggettマスター方程式(5.57)の導出に入った仮定をまとめよう。
\begin{itemize}
\item Born--Markov近似が成り立つ（4.2.1節参照）。
\item 調和振動子の環境は、高振動数カットオフ$\Lambda$をもつOhmicなスペクトル密度$J(\omega)\propto\omega$で記述される。
\item 環境は十分高温であり、その熱エネルギーは系の固有振動数で定まるエネルギー尺度より、はるかに大きい（$k_{\mathrm B}T\gg\Omega$）。
\item 系の固有振動数は、環境の高振動数カットオフ$\Lambda$より、はるかに小さい（$\Omega\ll\Lambda$）。
\end{itemize}
環境監視と固有の運動の相互作用により、最小不確定性Gaussianである指針状態が出現することを見いだした。これらは位置と運動量の両方でよく局在し、位相空間の古典的な点を近似する（5.2.6節参照）。この状態は、環境による位置座標の直接監視と、位置座標・運動量座標について対称な自己Hamiltonianをもつ系の固有の運動の、組み合わさった影響に対して最も頑健な、妥協として生じる。固有の運動は、運動量の重ね合わせから空間的重ね合わせを作ることを通じ、運動量座標が環境から直接監視されなくても、運動量のデコヒーレンスをもたらす。

位置の「直接的」デコヒーレンスと、運動量の「間接的」デコヒーレンスの違いは、それぞれの時間尺度が明確に異なることにも反映される。運動量のデコヒーレンスは、系の固有の発展が定める力学的時間尺度で起こる。これに対し、位置のデコヒーレンスの時間尺度は、本質的にその力学的時間尺度に依存せず、通常ははるかに短い。

\subsubsectionlink*{スピン・ボソン模型}
この模型は、スピン$1/2$粒子で表す中心二準位系が、互いに相互作用しない調和振動子の環境に、双線形に結合するものである。

系のHamiltonianにトンネル項を含めない簡単化した模型と、環境のOhmicなスペクトル密度について、系の「上向きスピン」と「下向きスピン」の間の干渉が、環境温度$T$に線形に増加する率で、指数関数的に減衰することを見いだした（5.3.1節参照）。

また、一般のスピン・ボソン模型のBorn--Markovマスター方程式も導いた。その形は［(5.142)参照］
\chfiveeq{204}
である。

% 原文 p.240
\hypertarget{srcpage:240}{}\label{srcpage:240}
右辺の丸括弧の項は、環境によってシフトした自己Hamiltonian $\hat H'_{\mathrm S}$のもとでの発展を表し、第2項は、系の$\hat\sigma_z$固有基底でのデコヒーレンスに対応し、最後の二項は二準位系の崩壊を記述する。

\subsubsectionlink*{スピン環境模型}
$\hat\sigma_z$座標が、他のスピンの環境に線形に結合する、トンネル効果をもつ中心スピンの模型を用い、環境による異なる優先指針基底の選択を示した。また、弱結合極限では、スピン環境を振動子環境に写像できることを示した。すなわち、スピン環境に弱く結合した系の縮約された運動は、同じ系が、陽に温度依存のスペクトル密度
\chfiveeq{205}
で記述する同等の振動子環境に結合したものとして記述できる。$J(\omega)$はスピン環境のスペクトル密度である。したがって、振動子環境の場合、すなわち量子Brown運動模型とスピン・ボソン模型で得た先の表式をすべて引き継ぎ、有効スペクトル密度(5.205)で評価すればよい。スピン・スピン模型からスピン・ボソン模型への写像と、振動子・スピン模型から量子Brown運動模型への写像を論じた。

\subsubsectionlink*{マスター方程式は一般的であり、係数とスペクトル密度は個別的である}
標準模型は、かなり一般的なマスター方程式で記述され、ある仮定のもとでは、異なる標準模型の間でさえ、似た形式的構造を取ることがある。例えば5.4.2節では、弱結合極限において、スピン環境模型は、振動子環境模型のマスター方程式と形式的に同一の方程式で記述できるという意味で、後者へ写像できることを見た。実際、振動子環境の方程式は、弱結合の場合に普遍的である\cite{ref211}。

しかし、環境とその物理的性質の違いを担うのは、マスター方程式の項に掛かる係数の具体的な形である。そして、この係数は環境のスペクトル密度に依存する。例えば、標準的なOhmicスペクトル密度で評価した量子Brown運動のマスター方程式は、一般に、温度依存の「有効」スペクトル密度(5.205)で評価したマスター方程式の運動とは、大きく異なる縮約された運動を記述する。

% 原文 p.241
\hypertarget{srcpage:241}{}\label{srcpage:241}
係数の具体的な関数形も、重要な役割を果たす。例えば、量子Brown運動の扱いでは、主に陽な時間依存性のないMarkov型の係数［(5.26)参照］に注目した。しかし、代わりに、陽に時間依存する係数［(5.42)参照］を用いて、デコヒーレンスの短時間の運動を、より現実的に扱う方法も示した。さらに、係数の厳密な表式を用いれば、完全な非Markov動力学が得られる（5.2.7節参照）。どの場合も、マスター方程式は同じ形(5.202)を取るが、同じスペクトル密度を用いても、係数の選び方の違いで、得られる運動は大きく異なりうる。

したがって、デコヒーレンス模型のマスター方程式の実際の物理は、主に、その評価に用いる係数とスペクトル密度の、具体的な形に含まれると結論する。

\subsubsectionlink*{一般性}
いうまでもなく、模型は物理系の簡単化した表現にすぎない。いくつかの模型に入った、重要な仮定と近似を論じた。
\begin{itemize}
\item Born--Markov近似。
\item 系と環境の双線形結合の仮定（ここで論じたすべての模型で採用した）。
\item 環境の暗黙の長波長仮定。これは、量子Brown運動における、デコヒーレンス率が飽和しないという非現実的な性質をもたらす。
\item マスター方程式の時間依存係数を、長時間の漸近値で近似すること［(5.26)、(5.42)を比較］。
\item Caldeira--Leggettマスター方程式の導出における、高温環境の仮定。
\end{itemize}
これらの仮定は、ある物理系の模型化ではまったく無害でも、別の系では非現実的な結果をもたらすことがある。したがって、具体的な実験状況を、特定のデコヒーレンスの標準模型で模型化する際には、各模型の基礎となる仮定と、対象となる実験状況でそれらがどの程度成り立つかを、正しく理解するよう、特別の注意を払わなければならない。
